\begin{lemma}{Erwartungswert des Supremums integrierbarer Funktionen}
            {EndlicheDarstellungDesSupremumsIntegrierbarerFunktionen}
Seien $I$ eine nichtleere, abzählbare Menge und 
$\folge X i I \in L_1(\Omega, \mfa, \Prim, \R)^I$.\\
Es gelte:
\begin{enumerate}[(a)]
\item für alle 
     $\omega \in \Omega$ ist $\menge{X_i(\omega)}{i \in I}$
     nach oben beschränkt,
\item $\Sup{i \in I} X_i \in L_1(\Omega, \mfa, \Prim, \R)$
(wohldefiniert wegen (a) und $I\neq \emptyset$).
\end{enumerate}
Dann gilt:
\begin{enumerate}[(1)]
\item\label{EndlicheDarstellungDesSupremumsIntegrierbarerFunktionen1}
 $\menge{\E\left( \Sup{i \in E} X_i \right)}{E \in \Pot_{ne}(I)} $
    ist nichtleer und nach oben beschränkt.
\item\label{EndlicheDarstellungDesSupremumsIntegrierbarerFunktionen2}
 $\sup\menge{\E\left( \Sup{i \in E} X_i\right)}{E \in \Pot_{ne}(I)} 
    = \E\left(\Sup{i \in I} X_i \right)$.
\end{enumerate}
\end{lemma}
\begin{proof}
Da $I$ nichtleer ist, gibt es ein $i_0 \in I$. Daher ist
$\meng{i_0}\in\Pot_{ne}(I)$, sodass die in (1) angegebene Menge nichtleer
ist.\\
Sei $A \in \Pot_{ne}(I)$. Dann ist auch $A \neq \emptyset$, sodass ein
$a \in A$ existiert.\\
Wegen
\[X_a \le \sup_{i \in A}X_i \le \sup_{i \in I}X_i =: f,\]
der Integrierbarkeit von $X_a$ und nach (b) auch der von $f$
ist nach \refclaim{Integrierbarkeitskriterien}{Integrierbarkeitskriterien2}
\begin{align}\sup_{i \in A}X_i \text{ integrierbar mit }
\E\left(\sup_{i \in A}X_i \right) \le  \E\left(f \right)\text. \tag{$\ast$}
\end{align}
Somit ist die in (1) angegebene Menge auch nach oben beschränkt.\\
Falls nun $I$ endlich ist, so ist wegen $I \in \Pot_{ne}(I)$ für (2) nichts zu
zeigen.\\
Sei also $I$ abzählbar unendlich.
Dann existiert $\iota: \N \rightarrow I$ bijektiv.\\
Setze
\[I_n \coloneqq  \iota(\haken n) \text{ und } f_n\coloneqq  \sup_{i \in I_n}X_i
\text{ für alle } n \in \N\text.\]
Dann ist $\folge I n \N$ eine aufsteigende Mengenfolge mit
$\bigcup\menge{I_n}{n \in \N} = I$ und daher $\folge f n \N$ eine monoton
wachsende Folge nach $(\ast)$ integrierbarer Funktionen, die punktweise gegen
$f$ konvergiert.\\
Somit sind die Voraussetzungen von \ref{MonotoneKonvergenzImIntegrierbarenFall}
erfüllt und es gilt:
\begin{align*}
\E(f)& \overset{\ref{MonotoneKonvergenzImIntegrierbarenFall}}= 
\Limes n \infty \E(f_n) = \sup_{n \in \N}\E(f_n)
= \sup\menge{\E\left(\sup_{i \in I_n}X_i \right)}{n \in \N}\\
&\le\sup\menge{\E\left(\sup_{i \in E}X_i \right)}{E \in \Pot_{ne}(I)}
\overset{(\ast)} \le \E(f)\text.
\end{align*}
\end{proof}
